
\documentclass[12pt]{article}

\usepackage[spanish]{babel}
\usepackage[utf8]{inputenc}
\usepackage[T1]{fontenc}
\usepackage{geometry}
\usepackage{graphicx}
\usepackage{hyperref}
\usepackage{amsmath}
\usepackage{array}
\usepackage{booktabs}

\geometry{a4paper, margin=2.5cm}

\title{Informe de Programación \ Gestión de Estudiantes y Notas}
\author{Nombre del Estudiante}
\date{\today}

\begin{document}

\maketitle

\tableofcontents
\newpage

\section{Formulación del Problema}

\subsection{Descripción del Problema}
Una escuela desea implementar un programa en C que le permita
gestionar las calificaciones de un grupo de estudiantes en varias asignaturas. Además de gestionar
las calificaciones, el programa debe calcular y mostrar información relevante como el promedio
de calificaciones por estudiante y por asignatura, la calificación más alta y baja, y cuántos
estudiantes aprobaron cada asignatura.

\subsection{Abstracción del Problema}

\subsubsection{Organización de los Datos}
Los datos se organizan mediante una matriz de dimensiones $5 \times 3$, donde:
\begin{itemize}
\item Filas: representan estudiantes.
\item Columnas: representan materias.
\end{itemize}

\begin{table}[h!]
\centering
\begin{tabular}{c|c|c|c}
\textbf{} & \textbf{Materia 1} & \textbf{Materia 2} & \textbf{Materia 3} \\ \hline
\textbf{Estudiante 1} & ? & ? & ? \\
\textbf{Estudiante 2} & ? & ? & ? \\
\textbf{Estudiante 3} & ? & ? & ? \\
\textbf{Estudiante 4} & ? & ? & ? \\
\textbf{Estudiante 5} & ? & ? & ? \\
\end{tabular}
\caption{Representación de las notas de los estudiantes}
\end{table}

\subsubsection{Cálculos Realizados}

Para el desarrollo del programa se implementaron diversos cálculos estadísticos aplicados a una matriz de dimensiones $5 \times 3$, donde cada fila representa un estudiante y cada columna una asignatura.

A continuación, se describen las fórmulas utilizadas:

\paragraph{Promedio por Estudiante}

El promedio de cada estudiante se calcula como la suma de sus calificaciones dividida para el número total de asignaturas:

[
P_i = \frac{n_{i1} + n_{i2} + n_{i3}}{3}
]

donde $P_i$ representa el promedio del estudiante $i$, y $n_{ij}$ corresponde a la nota del estudiante $i$ en la asignatura $j$.

\paragraph{Promedio por Asignatura}

El promedio de cada asignatura se obtiene sumando las calificaciones de todos los estudiantes en dicha materia y dividiéndolas para el número total de estudiantes:

[
P_j = \frac{n_{1j} + n_{2j} + n_{3j} + n_{4j} + n_{5j}}{5}
]

donde $P_j$ representa el promedio de la asignatura $j$.

\paragraph{Conteo de Aprobados y Reprobados}

Se establece como criterio de aprobación una nota mayor o igual a 6. Para el conteo de estudiantes aprobados por asignatura se utiliza la siguiente expresión:

[
A_j = \sum_{i=1}^{5}
\begin{cases}
1 & \text{si } n_{ij} \geq 6 \
0 & \text{si } n_{ij} < 6
\end{cases}
]

De forma complementaria, el número de estudiantes reprobados se calcula como:

[
R_j = \sum_{i=1}^{5}
\begin{cases}
1 & \text{si } n_{ij} < 6 \
0 & \text{si } n_{ij} \geq 6
\end{cases}
]

\subsubsection{Ejemplo Práctico}

Supongamos las notas de la \textbf{Materia 1} (\(j=1\)):

\[
\begin{array}{c|c|c|c}
\text{Estudiante} & \text{Nota} & \text{¿} \geq 6\text{?} & \text{Aporta} \\
\hline
i=1 & 7 & \text{Sí} & +1 \\
i=2 & 5 & \text{No} & +0 \\
i=3 & 8 & \text{Sí} & +1 \\
i=4 & 4 & \text{No} & +0 \\
i=5 & 9 & \text{Sí} & +1 \\
\end{array}
\]

\[
A_1 = 1 + 0 + 1 + 0 + 1 = 3 \text{ estudiantes aprobados}
\]


\paragraph{Implementación}

Las fórmulas descritas fueron implementadas en el programa mediante estructuras iterativas (\texttt{for}) para recorrer la matriz, y funciones auxiliares para determinar valores máximos y mínimos, garantizando así una solución modular y eficiente.

\subsubsection{Validación de Entradas}
Las notas deben estar dentro del rango permitido (0 a 10). En este caso, los valores son generados aleatoriamente, asegurando que cumplen con el rango establecido.

\subsubsection{Visualización de Resultados}
Los resultados se presentan en consola de forma estructurada, mostrando:
\begin{itemize}
\item Notas por estudiante.
\item Promedio individual.
\item Nota máxima y mínima.
\item Estadísticas por asignatura.
\end{itemize}

\section{Implementación del Código}

\subsection{Repositorio}
El código fuente se encuentra disponible en el siguiente enlace:
\begin{center}
\url{https://github.com/tu-usuario/tu-repositorio}
\end{center}

\subsection{Historial de Commits}
Se realizó un seguimiento del desarrollo mediante commits en GitHub, donde cada cambio fue documentado con una descripción clara de su propósito.

\begin{itemize}
\item Commit 1: Estructura inicial del programa.
\item Commit 2: Implementación de la matriz y carga de datos.
\item Commit 3: Cálculo de promedios por estudiante.
\item Commit 4: Cálculo de estadísticas por asignatura.
\item Commit 5: Mejora de la salida en consola.
\end{itemize}

\section{Validación y Pruebas}

Se realizaron múltiples ejecuciones del programa para verificar su correcto funcionamiento.

\subsection{Evidencia de Ejecución}

\begin{figure}[h!]
\centering
\includegraphics[width=0.8\textwidth]{ejecucion.png}
\caption{Ejemplo de ejecución del programa}
\end{figure}

\section{Conclusiones}

\subsection{Aprendizajes}
Durante el desarrollo de este programa se reforzaron los siguientes conceptos:
\begin{itemize}
\item Uso de matrices en lenguaje C.
\item Estructuras de control (bucles y condicionales).
\item Modularización mediante funciones.
\item Cálculo de estadísticas básicas.
\end{itemize}

\subsection{Retos Enfrentados}
Antes de comenzar a programar, los principales retos fueron:
\begin{itemize}
\item Comprender cómo estructurar los datos correctamente.
\item Definir los cálculos necesarios.
\item Evitar errores en los recorridos de la matriz.
\end{itemize}

\end{document}
